% !TEX program = lualatex
% Instructivo de instalacion del .exe - DUOC FPY1101
% Compilar: lualatex Instructivo-Instalacion.tex (dos veces para TOC)
\DocumentMetadata{
  pdfstandard   = ua-2,
  pdfversion    = 2.0,
  lang          = es-CL,
  tagging       = on,
  testphase     = {phase-III,tagpdf,table,firstaid}
}

\documentclass[11pt,a4paper,oneside]{report}

\usepackage{fontspec}
\usepackage{polyglossia}
\setdefaultlanguage{spanish}

\usepackage[a4paper,margin=2.5cm,top=2.8cm,bottom=2.8cm]{geometry}
\usepackage{xcolor}

% Colores contraste WCAG AA (>= 4.5:1 sobre blanco)
\definecolor{primary}{HTML}{0050A0}
\definecolor{accent}{HTML}{8B4500}
\definecolor{success}{HTML}{1B5E20}
\definecolor{danger}{HTML}{B71C1C}
\definecolor{boxBg}{HTML}{F4F6FA}
\definecolor{boxBorder}{HTML}{C5D0E0}
\definecolor{codeBg}{HTML}{F0F0F0}

\usepackage{hyperref}
\hypersetup{
  pdftitle    = {Instructivo de Instalacion - Entrega de Evaluacion},
  pdfauthor   = {DUOC UC - Roberto Arce - FPY1101},
  pdfsubject  = {Como instalar y ejecutar la aplicacion de entrega de evaluaciones},
  pdfkeywords = {instalacion, Windows, SmartScreen, exe, DUOC, FPY1101},
  pdflang     = {es-CL},
  colorlinks  = true, linkcolor = primary, urlcolor = primary,
  citecolor   = primary, bookmarksnumbered = true
}

\usepackage{titlesec}
\titleformat{\chapter}[display]
  {\normalfont\Huge\bfseries\color{primary}}
  {\filright\Large\MakeUppercase{\chaptertitlename}\ \thechapter}
  {1ex}{\filright}
\titlespacing*{\chapter}{0pt}{-30pt}{20pt}
\titleformat{\section}{\normalfont\Large\bfseries\color{primary}}{\thesection}{1em}{}
\titleformat{\subsection}{\normalfont\large\bfseries\color{accent}}{\thesubsection}{1em}{}

\usepackage{enumitem}
\setlist{itemsep=2pt,parsep=2pt}

\usepackage{tcolorbox}
\tcbuselibrary{skins,breakable}
\newtcolorbox{notebox}[1][]{enhanced,breakable,colback=boxBg,colframe=boxBorder,boxrule=0.5pt,arc=3pt,left=8pt,right=8pt,top=6pt,bottom=6pt,title=#1,fonttitle=\bfseries\color{primary},coltitle=primary,colbacktitle=boxBg,attach boxed title to top left={xshift=8pt,yshift=-9pt},boxed title style={size=small,colback=white,colframe=white}}
\newtcolorbox{warnbox}[1][]{enhanced,breakable,colback=boxBg,colframe=danger,boxrule=1pt,arc=3pt,left=8pt,right=8pt,top=6pt,bottom=6pt,title=#1,fonttitle=\bfseries\color{white},coltitle=white,colbacktitle=danger,attach boxed title to top left={xshift=8pt,yshift=-9pt},boxed title style={size=small,colback=danger,colframe=danger}}
\newtcolorbox{successbox}[1][]{enhanced,breakable,colback=boxBg,colframe=success,boxrule=1pt,arc=3pt,left=8pt,right=8pt,top=6pt,bottom=6pt,title=#1,fonttitle=\bfseries\color{white},coltitle=white,colbacktitle=success,attach boxed title to top left={xshift=8pt,yshift=-9pt},boxed title style={size=small,colback=success,colframe=success}}

\usepackage{fancyhdr}
\pagestyle{fancy}
\fancyhf{}
\renewcommand{\headrulewidth}{0.4pt}
\fancyhead[L]{\small\color{primary}Instalacion}
\fancyhead[R]{\small\color{primary}Entrega de Evaluacion}
\fancyfoot[C]{\thepage}

\renewcommand{\contentsname}{Indice}
\renewcommand{\chaptername}{Capitulo}

\begin{document}
\sloppy

\begin{titlepage}
\centering
\vspace*{2cm}
{\color{primary}\Huge\bfseries Instructivo de Instalacion}\\[8pt]
{\color{accent}\LARGE Aplicacion de Entrega de Evaluaciones}\\[2.5cm]
\rule{0.7\textwidth}{0.8pt}\\[10pt]
{\large DUOC UC\,---\,FPY1101 Programacion en Python}\\[10pt]
\rule{0.7\textwidth}{0.8pt}\\[3cm]
{\large Profesor: Roberto Arce}\\[8pt]
{\large Secciones 001D \textbullet{} 002D \textbullet{} 003D}\\[3cm]
{\small \href{https://github.com/Robbyfuu/entrega-evaluacion-github/releases/latest}{Descarga oficial en GitHub Releases}}\\
\vfill
\end{titlepage}

\tableofcontents
\clearpage

% ============================================================
\chapter{Que vas a instalar}

La aplicacion \textbf{Entrega de Evaluacion} es un programa de Windows que te
permite subir tus evaluaciones de programacion a GitHub de forma automatica.

\begin{notebox}[Ventaja]
No necesitas instalar Git, GitHub CLI ni Python por separado para subir tu
codigo: la aplicacion trae todo lo necesario incorporado. Solo necesitas
Python para programar (IDLE).
\end{notebox}

\section{Requisitos}
\begin{itemize}
  \item Windows 10 o 11 (64 bits)
  \item Conexion a internet
  \item Cuenta de GitHub (gratuita)
  \item Python instalado (para programar en IDLE)
\end{itemize}

\section{Dos formas de instalar}

\begin{itemize}
  \item \textbf{Instalador} (\texttt{EntregaEvaluacion-Setup.exe}):
    recomendado. Crea acceso directo y permite inicio automatico.
  \item \textbf{Ejecutable directo} (\texttt{EntregaEvaluacion.exe}):
    sin instalar, doble-click y listo. Util si no tienes permisos.
\end{itemize}

% ============================================================
\chapter{Descarga}

\section{Paso 1: Abrir la pagina de descargas}

Abre en tu navegador el enlace oficial:

\begin{center}
\texttt{github.com/Robbyfuu/entrega-evaluacion-github/releases/latest}
\end{center}

\section{Paso 2: Descargar el instalador}

En la seccion \textbf{Assets} descarga:

\begin{center}
\textbf{EntregaEvaluacion-Setup.exe}
\end{center}

\begin{notebox}[Si tu navegador advierte de la descarga]
Chrome o Edge pueden decir que el archivo "se descarga con poca frecuencia"
o "puede ser peligroso". Es normal en programas nuevos sin firma comercial.

Haz clic en la flecha junto a la descarga y elige \textbf{Conservar} o
\textbf{Mantener de todos modos}.
\end{notebox}

% ============================================================
\chapter{Instalacion y la advertencia de Windows}

\section{Paso 1: Ejecutar el instalador}

Haz doble-click en \texttt{EntregaEvaluacion-Setup.exe}.

\section{Paso 2: "Windows protegio su PC" (IMPORTANTE)}

\begin{warnbox}[Esta pantalla azul es NORMAL]
Como la aplicacion es nueva y no tiene una firma comercial de pago, Windows
muestra una pantalla azul llamada \textbf{SmartScreen}:

\vspace{4pt}
\textit{"Windows protegio su PC. Microsoft Defender SmartScreen impidio el
inicio de una aplicacion no reconocida."}
\vspace{4pt}

NO significa que el programa tenga virus. Solo significa que es nuevo.
\end{warnbox}

\textbf{Para continuar:}

\begin{enumerate}
  \item Haz clic en el texto \textbf{"Mas informacion"}.
  \item Aparece un boton nuevo abajo: \textbf{"Ejecutar de todas formas"}.
  \item Haz clic en \textbf{"Ejecutar de todas formas"}.
\end{enumerate}

\begin{notebox}[Por que es seguro]
El codigo fuente completo de la aplicacion es publico y revisable en GitHub.
Puedes verlo en el mismo repositorio desde donde descargaste el programa.
\end{notebox}

\section{Paso 3: Completar la instalacion}

\begin{enumerate}
  \item Elige la carpeta de instalacion (deja la que viene por defecto).
  \item Marca \textbf{"Crear acceso directo en el Escritorio"}.
  \item Si tu profesor lo indica, marca \textbf{"Iniciar automaticamente al
    encender el PC"}.
  \item Clic en \textbf{Instalar} y luego \textbf{Finalizar}.
\end{enumerate}

\begin{successbox}[Listo]
La aplicacion queda instalada con un icono en tu Escritorio.
\end{successbox}

\section{Alternativa: ejecutable directo (sin instalar)}

Si no puedes instalar (sin permisos de administrador):

\begin{enumerate}
  \item Descarga \texttt{EntregaEvaluacion.exe} (no el Setup).
  \item Guardalo en una carpeta facil de encontrar (ej: Escritorio).
  \item Doble-click. Aparece la misma pantalla azul de SmartScreen:
    \textbf{Mas informacion} $\rightarrow$ \textbf{Ejecutar de todas formas}.
\end{enumerate}

% ============================================================
\chapter{Primer uso}

\section{Paso 1: Elegir tu seccion}

La primera vez, la aplicacion pregunta tu seccion. Elige \textbf{001D},
\textbf{002D} o \textbf{003D} segun tu inscripcion.

\section{Paso 2: Iniciar sesion en GitHub}

\begin{enumerate}
  \item Clic en \textbf{Iniciar sesion}.
  \item Aparece un codigo grande tipo \texttt{XXXX-XXXX}.
  \item Abre en tu navegador (o celular): \url{https://github.com/login/device}
  \item Ingresa el codigo, inicia sesion en GitHub y autoriza.
  \item La aplicacion detecta la autorizacion automaticamente.
\end{enumerate}

\begin{notebox}[Una sola vez]
La sesion queda guardada. No tendras que volver a iniciar sesion en ese PC.
\end{notebox}

\section{Paso 3: Aceptar tu tarea}

Si tu profesor asigno tareas por GitHub Classroom, veras un aviso en la
aplicacion. Haz clic, acepta la tarea en el navegador, y vuelve.

% ============================================================
\chapter{Problemas comunes}

\section{No veo el boton "Ejecutar de todas formas"}

El texto \textbf{"Mas informacion"} es un enlace pequeno en la pantalla azul.
Haz clic ahi primero; el boton aparece despues.

\section{El antivirus borro el archivo}

Algunos antivirus muy estrictos mueven el \texttt{.exe} a cuarentena.
Restauralo desde el antivirus o pide ayuda a tu profesor. La descarga oficial
de GitHub es segura.

\section{Dice que falta un componente de Windows}

La aplicacion es autocontenida y no deberia faltar nada. Si aparece un error,
descarga la version mas reciente del enlace oficial.

\section{No se abre nada al hacer doble-click}

Espera unos segundos (la primera apertura descomprime archivos internos).
Si sigue sin abrir, prueba el instalador en lugar del ejecutable directo.

\section{Quiero desinstalar}

Panel de Control $\rightarrow$ Programas $\rightarrow$ Entrega de Evaluacion
$\rightarrow$ Desinstalar. O borra la carpeta si usaste el ejecutable directo.

% ============================================================
\chapter{Resumen rapido}

\begin{enumerate}
  \item Descargar \texttt{EntregaEvaluacion-Setup.exe} del enlace oficial.
  \item Doble-click. En la pantalla azul: \textbf{Mas informacion}
    $\rightarrow$ \textbf{Ejecutar de todas formas}.
  \item Instalar (acceso directo en Escritorio).
  \item Abrir, elegir seccion, iniciar sesion con codigo.
  \item Aceptar tarea (si el profesor la asigno).
  \item Programar, subir, entregar el enlace en el AVA.
\end{enumerate}

\begin{successbox}[Recuerda]
La pantalla azul de Windows es normal en programas nuevos. \textbf{Mas
informacion} $\rightarrow$ \textbf{Ejecutar de todas formas}. No es virus.
\end{successbox}

\end{document}
